\begin{standout}[第二章]
    半线性结构\\
    Semi-Linear Structures
\end{standout}

\begin{frame}{大纲}
    \setbeamertemplate{section in toc}[sections numbered]
    % \begin{multicols}{2}
    \tableofcontents
    % \end{multicols}
\end{frame}

\begin{frame}{学习目标}
    \begin{itemize}
        \item 理解树与二叉树的基本术语、性质和存储方式
        \item 能手工追踪前序、中序、后序及层次遍历
        \item 能追踪二叉搜索树操作、计算平均查找长度并理解平衡性
        \item 理解Huffman前缀码的构造，并能按码树编码和解码
    \end{itemize}
\end{frame}

\begin{frame}
    \begin{table}
        \small
        \label{tab:linear_structures_disadvantages}
        \caption{线性结构的优势与不足}
        \begin{tabular}{rcc}
            \toprule
            & \textbf{顺序列表} & \textbf{链式列表} \\
            \midrule
            访问元素 & $O(1)$ & $O(n)$ \\
            指定位置增删 & $O(n)$ & $O(1)$ \\
            \bottomrule
        \end{tabular}
    \end{table}
    \pause
    \centering
    链表增删的$O(1)$不包含查找结点的成本；树可表示层次关系，
    其操作代价取决于具体结构。
\end{frame}

\section{基本术语}

\begin{frame}
    \frametitle{\insertsectionhead}
    \begin{block}{树(tree)与森林(forest)}
        \begin{itemize}
            \item \textbf{半线性(semi-linear)}结构一般指\textbf{树}
            \item 树由$n$个\footnote{$n\in\mathbb{Z}\cap[0,+\infty)$}\textbf{顶
                  点(vertex)}\footnote{又名\textbf{结点(node)}}与连接于其间的若
                  干条\textbf{边(edge)}组成
            \item \alert{空树}既无结点亦无边
            \item \alert{非空树}应满足如下条件
                \begin{itemize}
                    \item 有且仅有$1$个特定结点为\textbf{根(root)}结点
                    \item 除根结点外的其余结点被分为$d$个
                        \footnote{$d\in\mathbb{Z}\cap[0,+\infty)$}\alert{互不相交}
                        的\textbf{子树(subtree)}
                    \item 子树与根之间由边相连，但不形成\textbf{环(ring)}
                \end{itemize} 
            \item 子树亦为树，满足上述性质（递归定义）
            \item $m$棵\footnote{$m\in\mathbb{Z}\cap[0,+\infty)$}互不相交的树的
                  集合为\textbf{森林}
        \end{itemize}
    \end{block}
\end{frame}

\begin{frame}
    \frametitle{\insertsectionhead}
    \begin{figure}
        \centering
        \begin{subfigure}[b]{0.15\textwidth}
            \centering
            \begin{tikzpicture}[ >=stealth, thick, black!50, %
                    list item/.style={draw=gray, circle, thick} %
                ]
                \node [terminal] (item1) at (0,0) {根};
            \end{tikzpicture}
            \caption{根树}
            \label{fig:tree_root_only}
        \end{subfigure}
        ~
        \begin{subfigure}[b]{0.2\textwidth}
            \centering
            \begin{tikzpicture}[ >=stealth, thick, black!50, %
                    list item/.style={draw=gray, circle, thick} %
                ]
                \node [terminal] (item1) at (0,0) {根};
                \node [terminal] (item2) at (-0.5,-1) {$\cdot$};
                \node [terminal] (item3) at (-1,-2) {$\cdot$};
                \node [terminal] (item4) at (-1.5,-3) {$\cdot$};
                \draw (item1) -- (item2) -- (item3) -- (item4);
            \end{tikzpicture}
            \caption{左斜树}
            \label{fig:tree_left_slop}
        \end{subfigure}
        ~
        \begin{subfigure}[b]{0.3\textwidth}
            \centering
            \begin{tikzpicture}[ >=stealth, thick, black!50, %
                    list item/.style={draw=gray, circle, thick} %
                ]
                \node [terminal] (item1) at (0,0) {根};
                \node [terminal] (item2) at (-0.5,-1) {$\cdot$};
                \node [terminal] (item3) at (0.5,-1) {$\cdot$};
                \node [terminal] (item4) at (-1,-2) {$\cdot$};
                \node [terminal] (item5) at (0,-2) {$\cdot$};
                \node [terminal] (item6) at (1,-2) {$\cdot$};
                \node [terminal] (item7) at (-1.5,-3) {$\cdot$};
                \node [terminal] (item8) at (-0.5,-3) {$\cdot$};
                \node [terminal] (item9) at (0.5,-3) {$\cdot$};
                \draw (item1) -- (item2);
                \draw (item1) -- (item3);
                \draw (item2) -- (item4);
                \draw (item2) -- (item5);
                \draw (item2) -- (item6);
                \draw (item4) -- (item7);
                \draw (item4) -- (item8);
                \draw (item6) -- (item9);
            \end{tikzpicture}
            \caption{普通树}
            \label{fig:tree_common}
        \end{subfigure}
        ~
        \begin{subfigure}[b]{0.3\textwidth}
            \centering
            \begin{tikzpicture}[ >=stealth, thick, black!50, %
                    list item/.style={draw=gray, circle, thick} %
                ]
                \node [terminal] (item1) at (0,0) {根};
                \node [terminal] (item2) at (-0.5,-1) {$\cdot$};
                \node [terminal] (item3) at (0.5,-1) {$\cdot$};
                \node [terminal] (item4) at (-1,-2) {$\cdot$};
                \node [terminal] (item5) at (0,-2) {$\cdot$};
                \node [terminal] (item6) at (1,-2) {$\cdot$};
                \node [terminal] (item7) at (-1.5,-3) {$\cdot$};
                \node [terminal] (item8) at (-0.5,-3) {$\cdot$};
                \node [terminal] (item9) at (0.5,-3) {$\cdot$};
                \node [terminal] (item10) at (1.5,-3) {$\cdot$};
                \draw (item1) -- (item2);
                \draw (item1) -- (item3);
                \draw (item2) -- (item4);
                \draw (item2) -- (item5);
                \draw (item2) -- (item6);
                \draw (item4) -- (item7);
                \draw (item4) -- (item8);
                \draw (item5) -- (item8);
                \draw (item5) -- (item9);
                \draw (item6) -- (item9);
            \end{tikzpicture}
            \caption{非树}
            \label{fig:nontree}
        \end{subfigure}
        \caption{几种树与非树}
        \label{fig:demo_trees_non_trees}
    \end{figure}
\end{frame}

\begin{frame}
    % \frametitle{\insertsectionhead}
    \begin{figure}
        \centering
        \begin{tikzpicture}[
            small mindmap,
            every node/.style={
                concept,
                },
            concept color=black!30!black!70,
            text width=4em,
            grow cyclic,
            align=flush center,
            level 1 concept/.append style={
                level distance=8em,
                font=\small,
                sibling angle=60,
                },
            level 2 concept/.append style={
                level distance=6em,
                font=\footnotesize,
                sibling angle=45
                },
        ]
            \node[
                root concept,
                font=\large
                ]{现代C++}[clockwise from=0]
            child[
                concept color={purple!30!black!70},
                level distance=12em,
                ]{node{过程式}[clockwise from=155]
                child{node{注释}}
                child{node{变量}}
                child{node{表达式}}
                child{node{引用}}
                child{node{函数}}
                child{node{程序\\结构}}
            }
            child[
                concept color=green!30!black!70,
                level distance=10em,
                ]{node{面向对象}[clockwise from=40]
                child{node{抽象\\与封装}}
                child{node{继承\\与多态}}
                child{node{运算符\\重载}}
                child{node{模板}[clockwise from=-30]}
            }
            child[
                concept color=red!30!black!70,
                level distance=10em,
                ]{node{函数式}[clockwise from=20]
                child{node{闭包}}
                child{node{纯函数}}
                child{node{惰性\\求值}}
                child{node{范围\\与单子}}
                child{node{元编程}}
            }
            child[
                concept color=blue!30!black!70,
                level distance=12em,
                ]{node{其他杂项}[clockwise from=-60]
                child{node{模块}}
                child{node{类型\\推断}}
                child{node{异常}}
                child{node{移动\\语义}}
                child{node{智能\\指针}}
                child{node{正则\\表达式}}
                child{node{并发}}
            };
        \end{tikzpicture}
        \caption{思维导图亦为树}
        \label{fig:mindmap_is_tree}
    \end{figure}
\end{frame}

\begin{frame}
    \frametitle{\insertsectionhead}
    \begin{block}{树的特点}
        \begin{itemize}
            \item 根结点无前驱，其余结点有且仅有$1$个直接前驱
            \item 所有结点均可有$n$个\footnote{$n\in\mathbb{Z}\cap[0,+\infty)$}
                  直接后继
            \item 前驱类似线性，后继则不同，故称半线性
        \end{itemize}
    \end{block}
\end{frame}

\begin{fragile}
    \frametitle{\insertsectionhead}
    \bicolumns[0.5]{
        \begin{block}{度(degree)}
            \begin{itemize}
                \item 结点的\textbf{度}指其子树个数
                \item 树的\textbf{度}指其最大结点度
                \item \textbf{叶(leaf)}结点度为$0$，亦称\textbf{终端}结点
                \item 其余结点为\textbf{分支}结点
            \end{itemize}
        \end{block}
    }{
        \begin{figure}
            \centering
            \begin{tikzpicture}[ >=stealth, thick, black!50, %
                    list item/.style={draw=gray, circle, thick}, %
                    level distance=5ex, sibling distance=5ex, %
                ]
                \node [terminal] {$2$}
                    child {node [terminal] {$3$}
                        child {node [terminal] {$2$}
                            child {node [terminal] {\alert{$0$}}}
                            child {node [terminal] {\alert{$0$}}}
                        }
                        child {node [terminal] {\alert{$0$}}}
                        child {node [terminal] {$1$}
                            child {node [terminal] {\alert{$0$}}}
                        }
                    }
                    child {node [terminal] {\alert{$0$}}};
            \end{tikzpicture}
            \caption{结点的度}
            \label{fig:degree_of_node}
        \end{figure}
    }
\end{fragile}

\begin{fragile}
    \frametitle{\insertsectionhead}
    \bicolumns[0.55]{
        \begin{block}{结点亲缘关系}
            \begin{itemize}
                \item 结点的子树根为该结点的\textbf{子(child)}结点
                      \[
                            b=child(a),\quad{}c=child(a)
                      \]
                \item 该结点为子树根的\textbf{父(parent)}结点
                      \[
                            a=parent(b),\quad{}a=parent(c)
                      \]
                \item 同一结点的所有子结点互为\textbf{兄弟(sibling)}
                      \[
                            b=sibling(c),\quad{}c=sibling(b)
                      \]
            \end{itemize}
        \end{block}
    }{
        \begin{figure}
            \centering
            \begin{tikzpicture}[ >=stealth, thick, black!50, %
                    list item/.style={draw=gray, circle, thick} %
                ]
                \node [terminal] {$a$}
                    child {node [terminal] {$b$}}
                    child {node [terminal] {$c$}};
            \end{tikzpicture}
            \caption{结点间的关系\footnotemark}
            \label{fig:relationship_of_node}
        \end{figure}
    }

    \footnotetext{$a$为$b$或$c$的父结点，$b$或$c$为$a$的子结点，$b$与$c$互为兄弟}
\end{fragile}

\begin{fragile}
    \frametitle{\insertsectionhead}
    \bicolumns[0.5]{
        \begin{block}{路径(path)、深度(depth)与高度(height)}
            \begin{itemize}
                \item 若结点序列$\{n_{i}\}_{i=0}^{k}$满足：
                      \[
                            n_{i}=parent(n_{i+1}),\quad{}i=0,1,\cdots,k-1
                      \]
                      则称之为自$n_{0}$至$n_{k}$的一条\textbf{路径}
                      \begin{itemize}
                        \item 所过\alert{边数}为路径\textbf{长度(length)}
                        \item 若存在自$n_{a}$至$n_{b}$的路径，则该路径\alert{唯
                              一}，且$n_{a}$为$n_{b}$的\textbf{祖先
                              (ancestor)}，$n_{b}$为$n_{a}$的\textbf{子孙
                              (descendant)}
                      \end{itemize}
                \item 结点\textbf{深度}\footnotemark{}为根至其的路径长度
                \item 结点\textbf{高度}为其最大子孙深度\footnotemark
                    \begin{itemize}
                        \item 树的高度为其根的高度
                    \end{itemize}
            \end{itemize}
        \end{block}
    }{
        \begin{figure}
            \centering
            \begin{tikzpicture}[ >=stealth, thick, black!50, %
                    list item/.style={draw=gray, circle, thick}, %
                    radius=2pt, scale=0.9,
                ]
                % \draw [help lines] (0,0) grid (5,6);
                \coordinate (ca) at (5.5,5.5);
                \coordinate (cb) at (5.5,3);
                \coordinate (cc) at (5.5,0.5);
                \coordinate (cd) at (-0.5,0);
                \coordinate (ce) at (-0.5,5.5);
                \coordinate (cr) at (2.5,5.5); % root
                \coordinate (cn) at (3.5,3); % n
                \coordinate (csa) at (3,2);
                \coordinate (csb) at (4,1.5);
                \coordinate (cta) at (2,4);
                \coordinate (ctb) at (1,3);
                \coordinate (ctc) at (1.5,2);
                \coordinate (ctd) at (0.5,0.5);
                \node [terminal] (root) at (cr) {根};
                \node [terminal] (n) at (cn) {$n$};
                \node at ($(cn)-(0,2)$) {\texttt{子树}};
                \draw [dashed,rounded corners=10pt]
                    (root) .. controls (0,4) .. (0,0)
                    -- (5,0) .. controls (5,4) .. (root);
                \draw [dashed,rounded corners=10pt]
                    (n) .. controls (2.5,2.5) .. (2.5,0.5)
                    -- ++(2,0) .. controls (4.5,2.5) .. (n);
                \draw [very thick,red!50!gray]
                    (root) .. controls ($(cr)+(0.5,-2)$) and ($(cn)+(-0.5,2)$) ..
                    node [sloped,above] {\smaller\texttt{\alert{路径}}} (n);
                \draw (n) .. controls ($(cn)+(-0.125,-0.75)$) and ($(csa)+(0.125,0.75)$) .. (csa) node [fill,circle] {};
                \draw (n) .. controls ($(cn)+(0.125,-1.5)$) and ($(csb)+(-0.125,1.25)$) .. (csb) node [fill,circle] {};
                \draw (root) .. controls ($(cr)+(-0.25,-1.25)$) and ($(cta)+(0,1.5)$) .. (cta) node [fill,circle] {}
                    .. controls ($(cta)+(-1.5,0)$) and ($(ctb)+(0,0.5)$) .. (ctb) node [fill,circle] {};
                \draw (cta) .. controls ($(cta)+(2,-1)$) and ($(ctc)+(0,1)$) .. (ctc) node [fill,circle] {}
                    .. controls ($(ctc)+(-1,0)$) and ($(ctd)+(0.5,1)$) .. (ctd) node [fill,circle] {};
                \draw [<->] (ca) -- node [sloped,above] {\texttt{结点\alert{深度}}} (cb);
                \draw [<->] (cb) -- node [sloped,above] {\texttt{结点\alert{高度}}} (cc);
                \draw [<->] (cd) -- node [sloped,above] {\texttt{树的\alert{高度}}} (ce);
                \draw (root.east) -- ($(ca)+(0.2,0)$);
                \draw (root.west) -- ($(ce)-(0.2,0)$);
                \draw (n.east) -- ($(cb)+(0.2,0)$);
                \draw ($(cc)-(1.2,0)$) -- ($(cc)+(0.2,0)$);
                \draw ($(cd)-(0.2,0)$) -- ($(cd)+(0.7,0)$);
            \end{tikzpicture}
            \caption{路径、高度与深度}
            \label{fig:path_depth_height}
        \end{figure}
    }

    \footnotetext[7]{又称结点所在\textbf{层数(layer)}，根在第$0$层}
    \footnotetext{此处范围仅限于以其为根的子树内，一般为该子树最大叶深}
\end{fragile}

\begin{frame}
    \frametitle{\insertsectionhead}
    \begin{block}{与线性结构的比较}
        \begin{table}
            \begin{tabular}{rl}
                \toprule
                \textbf{线性} & \textbf{半线性} \\
                \midrule
                首元素无前驱 & 根结点无父结点 \\
                尾元素无后继 & 叶结点无子结点 \\
                其他元素单前驱单后继 & 其他结点单父结点\alert{多}子结点 \\
                \bottomrule
            \end{tabular}
        \end{table}
    \end{block}
\end{frame}

\section{树}

% \begin{fragile}
%     \frametitle{\insertsectionhead}
%     \begin{block}{树的抽象数据类型}
%         \begin{minted}[linenos=false,escapeinside=@@]{console}
%             ADT Tree {
%             数据:
%                 数据对象: @$\mathcal{D} = \{a_{k} | a_{k}\in\text{数据元素集合}, k\in\mathbb{Z}\cap[1,n]\}$@
%                 逻辑关系: @$\mathcal{R} = \{\langle{}a_{k-1},a_{k}\rangle | k\in\mathbb{Z}\cap[2,n]\}$@
%             操作:
%                 create_stack()
%                     构造并初始化一个空栈@$s$@
%                 is_empty_stack()
%                     判断栈是否为空
%                 push_stack(s, x)
%                     若栈@$s$@存在且未满, 则在其栈顶插入值为@$x$@的新元素
%                 pop_stack(s)
%                     若栈@$s$@存在且非空, 则删除其栈顶元素
%                 top_stack(s)
%                     若栈@$s$@存在且非空, 则返回其栈顶元素
%             }
%         \end{minted}
%     \end{block}
% \end{fragile}

\begin{frame}
    \frametitle{\insertsectionhead}
    \begin{block}{树的存储结构}
        \begin{itemize}
            \item 可采用\alert{顺序}或\alert{链式}存储结构
            \item 每结点须记录：数据信息、与其他结点的逻辑关系
        \end{itemize}
    \end{block}
    \pause
    \begin{block}{树的结点关系表示方法}
        \begin{itemize}
            \item 父结点表示法：只记录父结点信息
            \item 子结点表示法：只记录子结点信息
            \item 父子结点表示法：同时记录父子结点信息
            \item 长子兄弟表示法：同时记录第一个子结点与兄弟结点信息
        \end{itemize}
    \end{block}
\end{frame}

\begin{fragile}
    \frametitle{\insertsectionhead}
    \bicolumns[0.5]{
        \begin{block}{父结点表示法}
            \begin{itemize}
                \item 采用数组按层存储各结点
                \item 每结点包括数据信息与父结点序号
            \end{itemize}
        \end{block}
        \begin{exampleblock}{复杂度}
            \begin{itemize}
                \item 空间：$O(n)$
                \item 时间：
                    \begin{itemize}
                        \item 查找父结点$O(1)$
                        \item 但查找子结点\alert{$O(n)$}
                    \end{itemize}
            \end{itemize}
        \end{exampleblock}
    }{
        \begin{figure}
            \centering
            \begin{bytefield}[boxformatting=\baselinealign]{2}
                \bitbox[]{2}{\tiny{$i$}}
                \bitbox{3}{\mintinline[fontsize=\smaller]{c}{data}}
                \bitbox{4}{\mintinline[fontsize=\smaller]{c}{parent}}
            \end{bytefield}
            \caption{父结点表示法}
            \label{fig:parent_representation}
        \end{figure}
        \begin{minted}{c}
            typedef struct {
                DataType data; // 数据信息
                int parent; // 父结点序号
            } TreeNode;
        \end{minted}
    }
\end{fragile}

\begin{fragile}
    \frametitle{\insertsectionhead}
    \begin{figure}
        \centering
        \begin{subfigure}[b]{0.4\textwidth}
            \centering
            \begin{tikzpicture}[ >=stealth, thick, black!50, %
                    list item/.style={draw=gray, circle, thick}, %
                    level distance=5ex, sibling distance=5ex, %
                ]
                \node [terminal] {$A$}
                    child {node [terminal] {$B$}
                        child {node [terminal] {$D$}
                            child {node [terminal] {$G$}}
                            child {node [terminal] {$H$}}
                        }
                        child {node [terminal] {$E$}}
                        child {node [terminal] {$F$}
                            child {node [terminal] {$I$}}
                        }
                    }
                    child {node [terminal] {$C$}};
            \end{tikzpicture}
            \caption{逻辑表示}
            \label{subfig:parent_logistic_representation}
        \end{subfigure}
        ~
        \begin{subfigure}[b]{0.4\textwidth}
            \centering
            \begin{bytefield}{2}
                \bitbox[]{2}{\tiny{$0$}}\bitboxes{3}{{$A$}{\alert{$-1$}}} \\
                \bitbox[]{2}{\tiny{$1$}}\bitboxes{3}{{$B$}{$0$}} \\
                \bitbox[]{2}{\tiny{$2$}}\bitboxes{3}{{$C$}{$0$}} \\
                \bitbox[]{2}{\tiny{$3$}}\bitboxes{3}{{$D$}{$1$}} \\
                \bitbox[]{2}{\tiny{$4$}}\bitboxes{3}{{$E$}{$1$}} \\
                \bitbox[]{2}{\tiny{$5$}}\bitboxes{3}{{$F$}{$1$}} \\
                \bitbox[]{2}{\tiny{$6$}}\bitboxes{3}{{$G$}{$3$}} \\
                \bitbox[]{2}{\tiny{$7$}}\bitboxes{3}{{$H$}{$3$}} \\
                \bitbox[]{2}{\tiny{$8$}}\bitboxes{3}{{$I$}{$5$}}
            \end{bytefield}
            \caption{存储表示}
            \label{subfig:parent_storage_representation}
        \end{subfigure}
        \caption{父结点表示法}
        \label{fig:demo_parent_representation}
    \end{figure}
\end{fragile}

\begin{fragile}
    \frametitle{\insertsectionhead}
    \bicolumns[0.5]{
        \begin{block}{子结点表示法}
            \begin{itemize}
                \item 采用数组按层存储各结点
                \item 每结点包括数据信息与子结点序号链表
            \end{itemize}
        \end{block}
        \begin{exampleblock}{复杂度}
            \begin{itemize}
                \item 空间：$O(n)$
                \item 时间：
                    \begin{itemize}
                        \item 查找子结点$O(d)$\footnotemark
                        \item 但查找父结点\alert{$O(n)$}
                    \end{itemize}
            \end{itemize}
        \end{exampleblock}
    }{
        \begin{figure}
            \centering
            \begin{bytefield}[boxformatting=\baselinealign]{2}
                \bitbox[]{2}{\tiny{$i$}}
                \bitbox{3}{\mintinline[fontsize=\smaller]{c}{data}}
                \bitbox{5}{\mintinline[fontsize=\smaller]{c}{children}}
            \end{bytefield}
            \caption{子结点表示法}
            \label{fig:children_representation}
        \end{figure}
        \begin{minted}{c}
            typedef struct {
                DataType data; // 数据信息
                LinkedList children; // 子结点序号链表
            } TreeNode;
        \end{minted}
    }

    \footnotetext{若该结点度数为$d$}
\end{fragile}

\begin{fragile}
    \frametitle{\insertsectionhead}
    \begin{figure}
        \centering
        \begin{subfigure}[b]{0.4\textwidth}
            \centering
            \begin{tikzpicture}[ >=stealth, thick, black!50, %
                    list item/.style={draw=gray, circle, thick}, %
                    level distance=5ex, sibling distance=5ex, %
                ]
                \node [terminal] {$A$}
                    child {node [terminal] {$B$}
                        child {node [terminal] {$D$}
                            child {node [terminal] {$G$}}
                            child {node [terminal] {$H$}}
                        }
                        child {node [terminal] {$E$}}
                        child {node [terminal] {$F$}
                            child {node [terminal] {$I$}}
                        }
                    }
                    child {node [terminal] {$C$}};
            \end{tikzpicture}
            \caption{逻辑表示}
            \label{subfig:children_logistic_representation}
        \end{subfigure}
        ~
        \begin{subfigure}[b]{0.4\textwidth}
            \centering
            \begin{bytefield}{2}
                \bitbox[]{2}{\tiny{$0$}}\bitbox{3}{$A$}\bitbox{2}{\tikzmark{saab}}\bitbox[]{2}{\tikzmark{taab}}\bitbox{3}{$1$}\bitbox{2}{\tikzmark{saac}}\bitbox[]{2}{\tikzmark{taac}}\bitbox{3}{$2$}\bitbox{2}{\tikzmark{saan}} \\
                \bitbox[]{2}{\tiny{$1$}}\bitbox{3}{$B$}\bitbox{2}{\tikzmark{sabd}}\bitbox[]{2}{\tikzmark{tabd}}\bitbox{3}{$3$}\bitbox{2}{\tikzmark{sabe}}\bitbox[]{2}{\tikzmark{tabe}}\bitbox{3}{$4$}\bitbox{2}{\tikzmark{sabf}}\bitbox[]{2}{\tikzmark{tabf}}\bitbox{3}{$5$}\bitbox{2}{\tikzmark{sabn}} \\
                \bitbox[]{2}{\tiny{$2$}}\bitbox{3}{$C$}\bitbox{2}{\tikzmark{sacn}} \\
                \bitbox[]{2}{\tiny{$3$}}\bitbox{3}{$D$}\bitbox{2}{\tikzmark{sadg}}\bitbox[]{2}{\tikzmark{tadg}}\bitbox{3}{$6$}\bitbox{2}{\tikzmark{sadh}}\bitbox[]{2}{\tikzmark{tadh}}\bitbox{3}{$7$}\bitbox{2}{\tikzmark{sadn}} \\
                \bitbox[]{2}{\tiny{$4$}}\bitbox{3}{$E$}\bitbox{2}{\tikzmark{saen}} \\
                \bitbox[]{2}{\tiny{$5$}}\bitbox{3}{$F$}\bitbox{2}{\tikzmark{safi}}\bitbox[]{2}{\tikzmark{tafi}}\bitbox{3}{$8$}\bitbox{2}{\tikzmark{safn}} \\
                \bitbox[]{2}{\tiny{$6$}}\bitbox{3}{$G$}\bitbox{2}{\tikzmark{sagn}} \\
                \bitbox[]{2}{\tiny{$7$}}\bitbox{3}{$H$}\bitbox{2}{\tikzmark{sahn}} \\
                \bitbox[]{2}{\tiny{$8$}}\bitbox{3}{$I$}\bitbox{2}{\tikzmark{sain}}
            \end{bytefield}
            \begin{tikzpicture}[ >=stealth, thick, black!50, overlay, remember picture, %
                    list item/.style={draw=gray, rounded corners, thick} %
                ]
                \draw[->] ($(pic cs:saab)+(0,-0.03)$) to [out=0,in=180] ($(pic cs:taab)+(0.26,-0.03)$);
                \fill ($(pic cs:saab)+(0,-0.03)$) circle [radius=0.05];
                \draw[->] ($(pic cs:saac)+(0,-0.03)$) to [out=0,in=180] ($(pic cs:taac)+(0.26,-0.03)$);
                \fill ($(pic cs:saac)+(0,-0.03)$) circle [radius=0.05];
                \draw[->] ($(pic cs:sabd)+(0,-0.03)$) to [out=0,in=180] ($(pic cs:tabd)+(0.26,-0.03)$);
                \fill ($(pic cs:sabd)+(0,-0.03)$) circle [radius=0.05];
                \draw[->] ($(pic cs:sabe)+(0,-0.03)$) to [out=0,in=180] ($(pic cs:tabe)+(0.26,-0.03)$);
                \fill ($(pic cs:sabe)+(0,-0.03)$) circle [radius=0.05];
                \draw[->] ($(pic cs:sabf)+(0,-0.03)$) to [out=0,in=180] ($(pic cs:tabf)+(0.26,-0.03)$);
                \fill ($(pic cs:sabf)+(0,-0.03)$) circle [radius=0.05];
                \draw[->] ($(pic cs:sadg)+(0,-0.03)$) to [out=0,in=180] ($(pic cs:tadg)+(0.26,-0.03)$);
                \fill ($(pic cs:sadg)+(0,-0.03)$) circle [radius=0.05];
                \draw[->] ($(pic cs:sadh)+(0,-0.03)$) to [out=0,in=180] ($(pic cs:tadh)+(0.26,-0.03)$);
                \fill ($(pic cs:sadh)+(0,-0.03)$) circle [radius=0.05];
                \draw[->] ($(pic cs:safi)+(0,-0.03)$) to [out=0,in=180] ($(pic cs:tafi)+(0.26,-0.03)$);
                \fill ($(pic cs:safi)+(0,-0.03)$) circle [radius=0.05];
                \draw ($(pic cs:saan)+(-0.23,-0.3)$) to ($(pic cs:saan)+(0.23,0.23)$);
                \draw ($(pic cs:sabn)+(-0.23,-0.3)$) to ($(pic cs:sabn)+(0.23,0.23)$);
                \draw ($(pic cs:sacn)+(-0.23,-0.3)$) to ($(pic cs:sacn)+(0.23,0.23)$);
                \draw ($(pic cs:sadn)+(-0.23,-0.3)$) to ($(pic cs:sadn)+(0.23,0.23)$);
                \draw ($(pic cs:saen)+(-0.23,-0.3)$) to ($(pic cs:saen)+(0.23,0.23)$);
                \draw ($(pic cs:safn)+(-0.23,-0.3)$) to ($(pic cs:safn)+(0.23,0.23)$);
                \draw ($(pic cs:sagn)+(-0.23,-0.3)$) to ($(pic cs:sagn)+(0.23,0.23)$);
                \draw ($(pic cs:sahn)+(-0.23,-0.3)$) to ($(pic cs:sahn)+(0.23,0.23)$);
                \draw ($(pic cs:sain)+(-0.23,-0.3)$) to ($(pic cs:sain)+(0.23,0.23)$);
            \end{tikzpicture}
            \vspace{-2ex}
            \caption{存储表示}
            \label{subfig:children_storage_representation}
        \end{subfigure}
        \caption{子结点表示法}
        \label{fig:demo_chidren_representation}
    \end{figure}
\end{fragile}

\begin{fragile}
    \frametitle{\insertsectionhead}
    \bicolumns[0.5]{
        \begin{block}{父子结点表示法}
            \begin{itemize}
                \item 结合上述二者
            \end{itemize}
        \end{block}
        \begin{exampleblock}{复杂度}
            \begin{itemize}
                \item 空间：$O(n)$
                \item 时间：
                    \begin{itemize}
                        \item 查找子结点$O(d)$
                        \item 但查找父结点$O(1)$
                    \end{itemize}
            \end{itemize}
        \end{exampleblock}
    }{
        \begin{figure}
            \centering
            \begin{bytefield}[boxformatting=\baselinealign]{2}
                \bitbox[]{2}{\tiny{$i$}}
                \bitbox{3}{\mintinline[fontsize=\smaller]{c}{data}}
                \bitbox{4}{\mintinline[fontsize=\smaller]{c}{parent}}
                \bitbox{5}{\mintinline[fontsize=\smaller]{c}{children}}
            \end{bytefield}
            \caption{父子结点表示法}
            \label{fig:parent_children_representation}
        \end{figure}
        \begin{minted}{c}
            typedef struct {
                DataType data; // 数据信息
                int parent; // 父结点序号
                LinkedList children; // 子结点序号链表
            } TreeNode;
        \end{minted}
    }
\end{fragile}

\begin{fragile}
    \frametitle{\insertsectionhead}
    \begin{figure}
        \centering
        \begin{subfigure}[b]{0.4\textwidth}
            \centering
            \begin{tikzpicture}[ >=stealth, thick, black!50, %
                    list item/.style={draw=gray, circle, thick}, %
                    level distance=5ex, sibling distance=5ex, %
                ]
                \node [terminal] {$A$}
                    child {node [terminal] {$B$}
                        child {node [terminal] {$D$}
                            child {node [terminal] {$G$}}
                            child {node [terminal] {$H$}}
                        }
                        child {node [terminal] {$E$}}
                        child {node [terminal] {$F$}
                            child {node [terminal] {$I$}}
                        }
                    }
                    child {node [terminal] {$C$}};
            \end{tikzpicture}
            \caption{逻辑表示}
            \label{subfig:parent_children_logistic_representation}
        \end{subfigure}
        ~
        \begin{subfigure}[b]{0.4\textwidth}
            \centering
            \begin{bytefield}{2}
                \bitbox[]{2}{\tiny{$0$}}\bitboxes{3}{{$A$}{\alert{$-1$}}}\bitbox{2}{\tikzmark{sbab}}\bitbox[]{2}{\tikzmark{tbab}}\bitbox{3}{$1$}\bitbox{2}{\tikzmark{sbac}}\bitbox[]{2}{\tikzmark{tbac}}\bitbox{3}{$2$}\bitbox{2}{\tikzmark{sban}} \\
                \bitbox[]{2}{\tiny{$1$}}\bitboxes{3}{{$B$}{$0$}}\bitbox{2}{\tikzmark{sbbd}}\bitbox[]{2}{\tikzmark{tbbd}}\bitbox{3}{$3$}\bitbox{2}{\tikzmark{sbbe}}\bitbox[]{2}{\tikzmark{tbbe}}\bitbox{3}{$4$}\bitbox{2}{\tikzmark{sbbf}}\bitbox[]{2}{\tikzmark{tbbf}}\bitbox{3}{$5$}\bitbox{2}{\tikzmark{sbbn}} \\
                \bitbox[]{2}{\tiny{$2$}}\bitboxes{3}{{$C$}{$0$}}\bitbox{2}{\tikzmark{sbcn}} \\
                \bitbox[]{2}{\tiny{$3$}}\bitboxes{3}{{$D$}{$1$}}\bitbox{2}{\tikzmark{sbdg}}\bitbox[]{2}{\tikzmark{tbdg}}\bitbox{3}{$6$}\bitbox{2}{\tikzmark{sbdh}}\bitbox[]{2}{\tikzmark{tbdh}}\bitbox{3}{$7$}\bitbox{2}{\tikzmark{sbdn}} \\
                \bitbox[]{2}{\tiny{$4$}}\bitboxes{3}{{$E$}{$1$}}\bitbox{2}{\tikzmark{sben}} \\
                \bitbox[]{2}{\tiny{$5$}}\bitboxes{3}{{$F$}{$1$}}\bitbox{2}{\tikzmark{sbfi}}\bitbox[]{2}{\tikzmark{tbfi}}\bitbox{3}{$8$}\bitbox{2}{\tikzmark{sbfn}} \\
                \bitbox[]{2}{\tiny{$6$}}\bitboxes{3}{{$G$}{$3$}}\bitbox{2}{\tikzmark{sbgn}} \\
                \bitbox[]{2}{\tiny{$7$}}\bitboxes{3}{{$H$}{$3$}}\bitbox{2}{\tikzmark{sbhn}} \\
                \bitbox[]{2}{\tiny{$8$}}\bitboxes{3}{{$I$}{$5$}}\bitbox{2}{\tikzmark{sbin}}
            \end{bytefield}
            \begin{tikzpicture}[ >=stealth, thick, black!50, overlay, remember picture, %
                    list item/.style={draw=gray, rounded corners, thick} %
                ]
                \draw[->] ($(pic cs:sbab)+(0,-0.03)$) to [out=0,in=180] ($(pic cs:tbab)+(0.26,-0.03)$);
                \fill ($(pic cs:sbab)+(0,-0.03)$) circle [radius=0.05];
                \draw[->] ($(pic cs:sbac)+(0,-0.03)$) to [out=0,in=180] ($(pic cs:tbac)+(0.26,-0.03)$);
                \fill ($(pic cs:sbac)+(0,-0.03)$) circle [radius=0.05];
                \draw[->] ($(pic cs:sbbd)+(0,-0.03)$) to [out=0,in=180] ($(pic cs:tbbd)+(0.26,-0.03)$);
                \fill ($(pic cs:sbbd)+(0,-0.03)$) circle [radius=0.05];
                \draw[->] ($(pic cs:sbbe)+(0,-0.03)$) to [out=0,in=180] ($(pic cs:tbbe)+(0.26,-0.03)$);
                \fill ($(pic cs:sbbe)+(0,-0.03)$) circle [radius=0.05];
                \draw[->] ($(pic cs:sbbf)+(0,-0.03)$) to [out=0,in=180] ($(pic cs:tbbf)+(0.26,-0.03)$);
                \fill ($(pic cs:sbbf)+(0,-0.03)$) circle [radius=0.05];
                \draw[->] ($(pic cs:sbdg)+(0,-0.03)$) to [out=0,in=180] ($(pic cs:tbdg)+(0.26,-0.03)$);
                \fill ($(pic cs:sbdg)+(0,-0.03)$) circle [radius=0.05];
                \draw[->] ($(pic cs:sbdh)+(0,-0.03)$) to [out=0,in=180] ($(pic cs:tbdh)+(0.26,-0.03)$);
                \fill ($(pic cs:sbdh)+(0,-0.03)$) circle [radius=0.05];
                \draw[->] ($(pic cs:sbfi)+(0,-0.03)$) to [out=0,in=180] ($(pic cs:tbfi)+(0.26,-0.03)$);
                \fill ($(pic cs:sbfi)+(0,-0.03)$) circle [radius=0.05];
                \draw ($(pic cs:sban)+(-0.23,-0.3)$) to ($(pic cs:sban)+(0.23,0.23)$);
                \draw ($(pic cs:sbbn)+(-0.23,-0.3)$) to ($(pic cs:sbbn)+(0.23,0.23)$);
                \draw ($(pic cs:sbcn)+(-0.23,-0.3)$) to ($(pic cs:sbcn)+(0.23,0.23)$);
                \draw ($(pic cs:sbdn)+(-0.23,-0.3)$) to ($(pic cs:sbdn)+(0.23,0.23)$);
                \draw ($(pic cs:sben)+(-0.23,-0.3)$) to ($(pic cs:sben)+(0.23,0.23)$);
                \draw ($(pic cs:sbfn)+(-0.23,-0.3)$) to ($(pic cs:sbfn)+(0.23,0.23)$);
                \draw ($(pic cs:sbgn)+(-0.23,-0.3)$) to ($(pic cs:sbgn)+(0.23,0.23)$);
                \draw ($(pic cs:sbhn)+(-0.23,-0.3)$) to ($(pic cs:sbhn)+(0.23,0.23)$);
                \draw ($(pic cs:sbin)+(-0.23,-0.3)$) to ($(pic cs:sbin)+(0.23,0.23)$);
            \end{tikzpicture}
            \vspace{-2ex}
            \caption{存储表示}
            \label{subfig:parent_children_storage_representation}
        \end{subfigure}
        \caption{父子结点表示法}
        \label{fig:demo_parent_chidren_representation}
    \end{figure}
\end{fragile}

\begin{frame}
    \frametitle{\insertsectionhead}
    \begin{alertblock}{父子结点表示法的性质}
        \begin{itemize}
            \item \textbf{优势}：一定程度上兼顾了查找效率
            \item \textbf{不足}：\alert{插入/删除}结点操作需大量修改链表，效率偏
                  低
        \end{itemize}
    \end{alertblock}
    \pause
    \begin{block}{基本术语}
        \begin{itemize}
            \item 若同一结点的所有子结点间具备某种线性次序，则称之为\textbf{有序
                  树(ordered tree)}
            \item 有序树的任意非叶结点均有且仅有$1$个\textbf{长子(eldest son)}
        \end{itemize}
    \end{block}
\end{frame}

\begin{fragile}
    \frametitle{\insertsectionhead}
    \bicolumns[0.5]{
        \begin{block}{长子兄弟表示法}
            \begin{itemize}
                \item 采用数组按层存储各结点
                \item 每结点包括
                \begin{itemize}
                    \item 数据信息
                    \item 长子结点序号
                    \item 首个兄弟结点序号
                \end{itemize}
                
            \end{itemize}
        \end{block}
    }{
        \begin{figure}
            \centering
            \begin{bytefield}[boxformatting=\baselinealign]{2}
                \bitbox[]{2}{\tiny{$i$}}
                \bitbox{3}{\mintinline[fontsize=\smaller]{c}{data}}
                \bitbox{7}{\mintinline[fontsize=\smaller]{c}{eldest_son}}
                \bitbox{5}{\mintinline[fontsize=\smaller]{c}{sibling}}
            \end{bytefield}
            \caption{长子兄弟表示法}
            \label{fig:eldest_son_sibling_representation}
        \end{figure}
        \begin{minted}{c}
            typedef struct {
                DataType data; // 数据信息
                int eldest_son; // 长子结点序号
                int sibling; // 兄弟结点序号
            } TreeNode;
        \end{minted}
    }
\end{fragile}

\begin{fragile}
    \frametitle{\insertsectionhead}
    \begin{figure}
        \centering
        \begin{subfigure}[b]{0.4\textwidth}
            \centering
            \begin{tikzpicture}[ >=stealth, thick, black!50, %
                    list item/.style={draw=gray, circle, thick}, %
                    level distance=5ex, sibling distance=5ex, %
                ]
                \node [terminal] {$A$}
                    child {node [terminal] {$B$}
                        child {node [terminal] {$D$}
                            child {node [terminal] {$G$}}
                            child {node [terminal] {$H$}}
                        }
                        child {node [terminal] {$E$}}
                        child {node [terminal] {$F$}
                            child {node [terminal] {$I$}}
                        }
                    }
                    child {node [terminal] {$C$}};
            \end{tikzpicture}
            \caption{逻辑表示}
            \label{subfig:eldest_son_sibling_logistic_representation}
        \end{subfigure}
        ~
        \begin{subfigure}[b]{0.4\textwidth}
            \centering
            \begin{bytefield}{2}
                \bitbox[]{2}{\tiny{$0$}}\bitboxes{3}{{$A$}{$1$}{\alert{$-1$}}} \\
                \bitbox[]{2}{\tiny{$1$}}\bitboxes{3}{{$B$}{$3$}{$2$}} \\
                \bitbox[]{2}{\tiny{$2$}}\bitboxes{3}{{$C$}{\alert{$-1$}}{\alert{$-1$}}} \\
                \bitbox[]{2}{\tiny{$3$}}\bitboxes{3}{{$D$}{$6$}{$4$}} \\
                \bitbox[]{2}{\tiny{$4$}}\bitboxes{3}{{$E$}{\alert{$-1$}}{$5$}} \\
                \bitbox[]{2}{\tiny{$5$}}\bitboxes{3}{{$F$}{$8$}{\alert{$-1$}}} \\
                \bitbox[]{2}{\tiny{$6$}}\bitboxes{3}{{$G$}{\alert{$-1$}}{$7$}} \\
                \bitbox[]{2}{\tiny{$7$}}\bitboxes{3}{{$H$}{\alert{$-1$}}{\alert{$-1$}}} \\
                \bitbox[]{2}{\tiny{$8$}}\bitboxes{3}{{$I$}{\alert{$-1$}}{\alert{$-1$}}}
            \end{bytefield}
            \caption{存储表示}
            \label{subfig:eldest_son_sibling_storage_representation}
        \end{subfigure}
        \caption{长子兄弟表示法}
        \label{fig:demo_eldest_son_sibling_representation}
    \end{figure}
\end{fragile}
\section{二叉树}

\begin{fragile}
    \frametitle{\insertsectionhead}
    \begin{block}{二叉树(binary tree)}
        \begin{itemize}
            \item 度不大于$2$的有序树
            \item 子结点可按左右区分
        \end{itemize}
    \end{block}
    \begin{alertblock}{将树转化为二叉树}
        \begin{itemize}
            \item 令长子为左子结点、首个兄弟为右子结点
            \item 任何树均可按此法转化为二叉树
            \item 因二叉树的表示与运算相对方便，故树的问题均可转化为二叉树形式进
                  行研究
        \end{itemize}
    \end{alertblock}
\end{fragile}

\begin{fragile}
    \frametitle{\insertsectionhead}
    \begin{figure}
        \centering
        \begin{subfigure}[b]{0.35\textwidth}
            \centering
            \begin{tikzpicture}[ >=stealth, thick, black!50, %
                    list item/.style={draw=gray, circle, thick}, %
                    level distance=5ex, % sibling distance=5ex, %
                    level/.style={sibling distance=12ex/#1}, %
                    % edge from parent path={
                    %     (\tikzparentnode.south) .. controls +(0,-0.2) and +(0,0.2) .. (\tikzchildnode.north)
                    % }, %
                ]
                \node [terminal] {$A$}
                    child {node [terminal] {$B$}
                        child {node [terminal] {$D$}
                            child {node [terminal] {$G$}}
                            child {node [terminal] {$H$}}
                        }
                        child {node [terminal] {$E$}}
                        child {node [terminal] {$F$}
                            child {node [terminal] {$I$}}
                        }
                    }
                    child {node [terminal] {$C$}};
            \end{tikzpicture}
            \caption{树表示}
            \label{subfig:tree_representation}
        \end{subfigure}
        ~
        \begin{subfigure}[b]{0.25\textwidth}
            \centering
            \begin{tikzpicture}[ >=stealth, thick, black!50, %
                    list item/.style={draw=gray, circle, thick}, %
                    % level distance=3ex, % sibling distance=5ex, %
                ]
                \matrix [row sep=2.5ex,column sep=3ex] {
                    \node[terminal](a){$A$}; & & \\
                    \node[terminal](b){$B$}; & \node[terminal](c){$C$}; & \\
                    \node[terminal](d){$D$}; & \node[terminal](e){$E$}; & \node[terminal](f){$F$}; \\
                    \node[terminal](g){$G$}; & \node[terminal](h){$H$}; & \node[terminal](i){$I$}; \\
                };
                \draw (a) -- (b) -- (d) -- (g);
                \draw (f) -- (i);
                \draw [dashed] (b) -- (c);
                \draw [dashed] (d) -- (e) -- (f);
                \draw [dashed] (g) -- (h);
            \end{tikzpicture}
            \vspace{-1ex}
            \caption{中间表示}
            \label{subfig:middle_representation}
        \end{subfigure}
        ~
        \begin{subfigure}[b]{0.35\textwidth}
            \centering
            \begin{tikzpicture}[ >=stealth, thick, black!50, %
                    list item/.style={draw=gray, circle, thick}, %
                    level distance=3.5ex, %
                    level/.style={sibling distance=16ex/#1}, %
                    % edge from parent path={
                    %     (\tikzparentnode.south) .. controls +(0,-0.1) and +(0,0.1) .. (\tikzchildnode.north)
                    % }, %
                ]
                \node [terminal] {$A$}
                    child {node [terminal] {$B$}
                        child {node [terminal] {$D$}
                            child {node [terminal] {$G$}
                                child [missing]
                                child {node [terminal] {$H$}}
                            }
                            child {node [terminal] {$E$}
                                child [missing]
                                child {node [terminal] {$F$}
                                    child {node [terminal] {$I$}}
                                    child [missing]
                                }
                            }
                        }
                        child {node [terminal] {$C$}}
                    }
                    child [missing];
            \end{tikzpicture}
            \caption{二叉树表示}
            \label{subfig:binary_tree_representation}
        \end{subfigure}
        \caption{将树转化为二叉树}
        \label{fig:demo_tree_to_binary_tree}
    \end{figure}
\end{fragile}

\subsection{几种特殊的二叉树}

\begin{fragile}
    \frametitle{\insertsubsectionhead}
    \bicolumns[0.5]{
        \begin{block}{左斜树\footnotemark}
            \begin{itemize}
                \item 所有非叶结点均有且仅有$1$个左子树
                \item 每层均有且仅有$1$个结点
                \item 已退化为线性结构
            \end{itemize}
        \end{block}
    }{
        \begin{figure}
            \centering
            \begin{subfigure}[b]{0.3\textwidth}
                \centering
                \begin{tikzpicture}[ >=stealth, thick, black!50, %
                        list item/.style={draw=gray, circle, thick}, %
                        level distance=4ex, sibling distance=4ex, %
                    ]
                    \node [terminal] {$A$}
                        child {node [terminal] {$B$}
                            child {node [terminal] {$C$}
                                child {node [terminal] {$D$}}
                                child [missing]
                            }
                            child [missing]
                        }
                        child [missing];
                \end{tikzpicture}
                \caption{左斜树}
                \label{subfig:left_slanted_tree}
            \end{subfigure}
            ~
            \begin{subfigure}[b]{0.3\textwidth}
                \centering
                \begin{tikzpicture}[ >=stealth, thick, black!50, %
                        list item/.style={draw=gray, circle, thick}, %
                        level distance=4ex, sibling distance=4ex, %
                    ]
                    \node [terminal] {$A$}
                        child [missing]
                        child {node [terminal] {$B$}
                            child [missing]
                            child {node [terminal] {$C$}
                                child [missing]
                                child {node [terminal] {$D$}}
                            }
                        };
                \end{tikzpicture}
                \caption{右斜树}
                \label{subfig:right_slanted_tree}
            \end{subfigure}
            \caption{斜树}
            \label{fig:demo_slanted_tree}
        \end{figure}
    }

    \footnotetext{右斜树与之类似，只需将左改作右}
\end{fragile}

\begin{fragile}
    \frametitle{\insertsubsectionhead}
    \bicolumns[0.5]{
        \begin{block}{满二叉树(full binary tree)}
            \begin{itemize}
                \item 所有叶结点均在最后一层
                \item 所有非叶结点度均为$2$
            \end{itemize}
        \end{block}
        \begin{exampleblock}{性质}
            \begin{itemize}
                \item 同深度的二叉树中满二叉树结点最多
                \item 同深度的二叉树中满二叉树叶结点最多
            \end{itemize}
        \end{exampleblock}
    }{
        \begin{figure}
            \centering
                \begin{tikzpicture}[ >=stealth, thick, black!50, %
                        list item/.style={draw=gray, circle, thick}, %
                        level distance=4ex, %
                        level 1/.style={sibling distance=16ex}, %
                        level 2/.style={sibling distance=8ex}, %
                        level 3/.style={sibling distance=4ex}, %
                    ]
                    \node [terminal] {$A$}
                        child {node [terminal] {$B$}
                            child {node [terminal] {$D$}
                                child {node [terminal] {$H$}}
                                child {node [terminal] {$I$}}
                            }
                            child {node [terminal] {$E$}
                                child {node [terminal] {$J$}}
                                child {node [terminal] {$K$}}
                            }
                        }
                        child {node [terminal] {$C$}
                            child {node [terminal] {$F$}
                                child {node [terminal] {$L$}}
                                child {node [terminal] {$M$}}
                            }
                            child {node [terminal] {$G$}
                                child {node [terminal] {$N$}}
                                child {node [terminal] {$O$}}
                            }
                        };
                \end{tikzpicture}
            \caption{满二叉树}
            \label{fig:demo_full_binary_tree}
        \end{figure}
    }
\end{fragile}

\begin{fragile}
    \frametitle{\insertsubsectionhead}
        \begin{figure}
            \centering
            \begin{subfigure}[b]{0.4\textwidth}
                \centering
                \begin{tikzpicture}[ >=stealth, thick, black!50, %
                        list item/.style={draw=gray, circle, thick}, %
                        level distance=4ex, %
                        level 1/.style={sibling distance=16ex}, %
                        level 2/.style={sibling distance=8ex}, %
                        level 3/.style={sibling distance=4ex}, %
                    ]
                    \node [terminal] {$A$}
                        child {node [terminal] {$B$}
                            child {node [terminal] {$D$}
                                child {node [terminal] {$H$}}
                                child {node [terminal] {$I$}}
                            }
                            child {node [terminal] {$E$}
                                child {node [terminal] {$J$}}
                                child {node [terminal] {$K$}}
                            }
                        }
                        child {node [terminal] {$C$}
                            child {node [terminal] {$F$}
                                child {node [terminal] {$L$}}
                                child {node [terminal] {$M$}}
                            }
                            child {node [terminal] {\alert{$G$}}}
                        };
                \end{tikzpicture}
                \caption{叶结点不在同层}
                \label{subfig:non_full_binary_tree_a}
            \end{subfigure}
            ~
            \begin{subfigure}[b]{0.4\textwidth}
                \centering
                \begin{tikzpicture}[ >=stealth, thick, black!50, %
                        list item/.style={draw=gray, circle, thick}, %
                        level distance=4ex, %
                        level 1/.style={sibling distance=16ex}, %
                        level 2/.style={sibling distance=8ex}, %
                        level 3/.style={sibling distance=4ex}, %
                    ]
                    \node [terminal] {$A$}
                        child {node [terminal] {$B$}
                            child {node [terminal] {$D$}
                                child {node [terminal] {$H$}}
                                child {node [terminal] {$I$}}
                            }
                            child {node [terminal] {$E$}
                                child {node [terminal] {$J$}}
                                child {node [terminal] {$K$}}
                            }
                        }
                        child {node [terminal] {$C$}
                            child {node [terminal] {$F$}
                                child {node [terminal] {$L$}}
                                child {node [terminal] {$M$}}
                            }
                            child {node [terminal] {\alert{$G$}}
                                child {node [terminal] {$N$}}
                                child [missing]
                            }
                        };
                \end{tikzpicture}
                \caption{个别非叶结点度不为$2$}
                \label{subfig:non_full_binary_tree_b}
            \end{subfigure}
            \caption{非满二叉树}
            \label{fig:demo_non_full_binary_tree}
        \end{figure}
\end{fragile}

\begin{fragile}
    \frametitle{\insertsubsectionhead}
    \bicolumns[0.5]{
        \begin{block}{完全二叉树(proper binary tree)}
            \begin{itemize}
                \item 若去除最后一层结点，则为满二叉树
                \item 最后一层结点自左至右连续\footnotemark{}排列
            \end{itemize}
        \end{block}
        \begin{exampleblock}{性质}
            \begin{itemize}
                \item 同结点数的二叉树中完全二叉树最矮
                \item 满二叉树亦为完全二叉树的一种
            \end{itemize}
        \end{exampleblock}
    }{
        \begin{figure}
            \centering
                \begin{tikzpicture}[ >=stealth, thick, black!50, %
                        list item/.style={draw=gray, circle, thick}, %
                        level distance=4ex, %
                        level 1/.style={sibling distance=16ex}, %
                        level 2/.style={sibling distance=8ex}, %
                        level 3/.style={sibling distance=4ex}, %
                    ]
                    \node [terminal] {$A$}
                        child {node [terminal] {$B$}
                            child {node [terminal] {$D$}
                                child {node [terminal] {$H$}}
                                child {node [terminal] {$I$}}
                            }
                            child {node [terminal] {$E$}
                                child {node [terminal] {$J$}}
                                child [missing]
                            }
                        }
                        child {node [terminal] {$C$}
                            child {node [terminal] {$F$}}
                            child {node [terminal] {$G$}}
                        };
                \end{tikzpicture}
            \caption{完全二叉树}
            \label{fig:demo_proper_binary_tree}
        \end{figure}
    }

    \footnotetext{指中间无空结点}
\end{fragile}
\subsection{二叉树的性质}

\begin{frame}
    \frametitle{\insertsubsectionhead}
    \begin{block}{性质甲}
        \begin{itemize}
            \item 令$N$层二叉树第$k$层结点数为$n_{l}(k)$，则有
            \[
                1\leq{}n_{l}(k)\leq2^{k},\;k\in\mathbb{Z}\cap[0,N)
            \]
        \end{itemize}
    \end{block}
    \pause
    \begin{exampleblock}{证明}
        \begin{enumerate}
            \item 左侧不等式显然成立
            \item 右侧不等式当$k=0$时，第$0$层仅有根结点，故$n_{l}(0)=1=2^{0}$显
                  然成立
            \item 假设当$k=n<N-1$时成立，即$n_{l}(n)\leq2^{n}$，因所有结点的度均
                  不大于$2$，故
                  \[
                        n_{l}(n+1)\leq2\cdot{}n_{l}(n)\leq2\cdot2^{n}=2^{n+1}
                  \]
                  于是当$k=n+1$时，归纳假设成立\hfill$\qed$
        \end{enumerate}
    \end{exampleblock}
\end{frame}

\begin{frame}
    \frametitle{\insertsubsectionhead}
    \begin{block}{性质乙}
        \begin{itemize}
            \item 令$k$层二叉树的结点总数为$n(k)$，则有
                  \[
                        k\leq{}n(k)\leq2^{k}-1
                  \]
        \end{itemize}
    \end{block}
    \pause
    \begin{exampleblock}{证明}
        \begin{enumerate}
            \item 由性质甲，$1\leq{}n_{l}(i)\leq2^{i},\;i\in\mathbb{Z}\cap[0,k)$
            \item 经累加后，有
                  \[
                        k=\sum_{i=0}^{k-1}{1}\leq{}n(k)=\sum_{i=0}^{k-1}{n_{l}(i)}\leq\sum_{i=0}^{k-1}{2^{i}}=2^{k}-1
                  \]
                  \hfill$\qed$
        \end{enumerate}
    \end{exampleblock}
\end{frame}

\begin{frame}
    \frametitle{\insertsubsectionhead}
    \begin{alertblock}{思考}
        \begin{itemize}
            \item 满足$n(k)=k$的二叉树一定是斜树么？
            \item 满足$n(k)=2^{k}-1$的二叉树一定是满二叉树么？
        \end{itemize}
    \end{alertblock}
\end{frame}

\begin{frame}
    \frametitle{\insertsubsectionhead}
    \begin{block}{性质丙}
        \begin{itemize}
            \item 令二叉树中度为$d$的结点数为$n_{d},\;(d\in\{0,1,2\})$，则有
                  \[
                        n_{0}=n_{2}+1
                  \]
        \end{itemize}
    \end{block}
    \pause
    \begin{exampleblock}{证明}
        \begin{enumerate}
            \item 一方面，结点总数由各种度的结点构成，故总结点数$n$满足
                  \[
                        n=\sum_{d=0}^{2}{n_{d}}=n_{0}+n_{1}+n_{2}
                  \]
            \item 另一方面，每个非根结点均由$1$个结点生成，故度为$d$的结点可生成
                  $d$个结点，即
                  \[
                        n=\sum_{d=0}^{2}{d\cdot{}n_{d}}+1=n_{1}+2n_{2}+1
                  \]
            \item 综上得证\hfill$\qed$
        \end{enumerate}
    \end{exampleblock}
\end{frame}

\begin{frame}
    \frametitle{\insertsubsectionhead}
    \begin{alertblock}{思考}
        \begin{itemize}
            \item 有$n$个结点的完全二叉树有多少叶结点？
        \end{itemize}
    \end{alertblock}
    \pause
    \begin{exampleblock}<2-|handout:0>{提示}
        \begin{itemize}
            \item 在完全二叉树中，度为$1$的结点数不多于$1$
            \item 当$n$为偶数时，$n_{1}=1$，$n_{0}=n/2$，$n_{2}=n/2-1$
            \item 当$n$为奇数时，$n_{1}=0$，$n_{0}=(n+1)/2$，$n_{2}=(n-1)/2$
            \item 故$n_{0}=\lceil{n/2}\rceil,\;n_{2}=\lfloor{(n-1)/2}\rfloor$
        \end{itemize}
    \end{exampleblock}
\end{frame}

\begin{frame}
    \frametitle{\insertsubsectionhead}
    \begin{block}{性质丁}
        \begin{itemize}
            \item 具有$n$个结点的完全二叉树层数$k=\lfloor\log_{2}{n}\rfloor+1$
        \end{itemize}
    \end{block}
    \pause
    \begin{exampleblock}{证明}
        \begin{enumerate}
            \item 由完全二叉树性质与性质乙可得，$k$层完全二叉树结点数$n$满足
                \[
                    2^{k-1}-1<n\leq2^{k}-1\;\text{或}\;2^{k-1}\leq{}n<2^{k}
                \]
            \item 取对数
                \[
                    k-1\leq\log_{2}{n}<k\;\text{或}\;\log_{2}{n}<k\leq\log_{2}{n}+1
                \]
            \item 注意到$k\in\mathbb{Z}$，于是得证\hfill$\qed$
        \end{enumerate}
    \end{exampleblock}
\end{frame}

\begin{fragile}
    \frametitle{\insertsubsectionhead}
    \bicolumns[0.5]{
        \begin{exampleblock}{完全二叉树按层序编号}
            \begin{itemize}
                \item 可为完全二叉树结点按层序依次编号
                \item 约定根编号为$1$，依次递增
                \item 称如此编号为$k$的结点为\textbf{结点$k$}
            \end{itemize}
        \end{exampleblock}
    }{
        \begin{figure}
            \centering
                \begin{tikzpicture}[ >=stealth, thick, black!50, %
                        list item/.style={draw=gray, circle, thick}, %
                        level distance=4ex, %
                        level 1/.style={sibling distance=16ex}, %
                        level 2/.style={sibling distance=8ex}, %
                        level 3/.style={sibling distance=4ex}, %
                    ]
                    \node [terminal] {$1$}
                        child {node [terminal] {$2$}
                            child {node [terminal] {$4$}
                                child {node [terminal] {$8$}}
                                child {node [terminal] {$9$}}
                            }
                            child {node [terminal] {$5$}
                                child {node [terminal] {$10$}}
                                child [missing]
                            }
                        }
                        child {node [terminal] {$3$}
                            child {node [terminal] {$6$}}
                            child {node [terminal] {$7$}}
                        };
                \end{tikzpicture}
            \caption{为完全二叉树按层序编号}
            \label{fig:demo_number_proper_binary_tree_by_layer_order}
        \end{figure}
    }
\end{fragile}

\begin{frame}
    \frametitle{\insertsubsectionhead}
    \begin{block}{性质戊}
        \begin{itemize}
            \item 为含有$n$个结点的完全二叉树按层序对结点编号\footnote{根结点为
                  $1$}，则有
                \begin{enumerate}
                    \item 结点$k$的左右子结点序号分别为$2k$与
                        $2k+1,\;k\in\mathbb{Z}\cap[1,\lfloor{}n/2\rfloor]$
                    \item 结点$k$的父结点序号为
                        $\lfloor{}k/2\rfloor,\;k\in\mathbb{Z}\cap(1,n]$
                \end{enumerate}
        \end{itemize}
    \end{block}
    \pause
    \begin{exampleblock}{证明}
        \begin{enumerate}
            \item 考察结论$1$，当$k=1$时，显然其左右子结点序号分别为$2$与$3$，成
                  立
            \item 假设当$k=m$时成立，即结点$m$的左右子结点序号分别为$2m$与
                  $2m+1$，
                  \begin{itemize}
                    \item 因结点$m+1$的左子结点必为结点$m$的右子结点的后继
                    \item 故结点$m+1$的左子结点序号为$(2m+1)+1=2(m+1)$
                    \item 且结点$m+1$的右子结点序号为$2(m+1)+1$
                  \end{itemize}
                  则当$k=m+1$时，假设成立，故结论$1$成立
            \item 由结论$1$知结论$2$成立\hfill\qed
        \end{enumerate}
    \end{exampleblock}
\end{frame}
\subsection{二叉树的存储结构}

\begin{frame}
    \frametitle{\insertsubsectionhead}
    \begin{block}{顺序存储}
        \begin{itemize}
            \item 按完全二叉树层序编号方式为二叉树编号，跳过不存在结点的编号
            \item 以静态数组方式存储，留空不存在的结点
        \end{itemize}
    \end{block}
\end{frame}

\begin{frame}
    \frametitle{\insertsubsectionhead}
    \begin{figure}
        \centering
        \begin{subfigure}[b]{0.45\textwidth}
            \centering
            \begin{tikzpicture}[ >=stealth, thick, black!50, %
                    list item/.style={draw=gray, circle, thick}, %
                    level distance=4ex, %
                    level 1/.style={sibling distance=16ex}, %
                    level 2/.style={sibling distance=8ex}, %
                    level 3/.style={sibling distance=4ex}, %
                ]
                \node [terminal] {$A$}
                    child {node [terminal] {$B$}
                        child {node [terminal] {$D$}}
                        child {node [terminal] {$E$}
                            child {node [terminal] {$G$}}
                            child [missing]
                        }
                    }
                    child {node [terminal] {$C$}
                        child [missing]
                        child {node [terminal] {$F$}}
                    };
            \end{tikzpicture}
            \caption{逻辑结构}
            \label{subfig:demo_logistic_structure_binary_tree_sequence_storage}
        \end{subfigure}
        ~
        \begin{subfigure}[b]{0.45\textwidth}
            \centering
            \begin{bytefield}{2}
                % \foreach \i in {0,...,10} {\bitbox[]{2}{\tiny{\i}}} \\
                \bitboxes[]{2}{{\tiny{$0$}}{\tiny{$1$}}{\tiny{$2$}}{\tiny{$3$}}{\tiny{$4$}}{\tiny{$5$}}{\tiny{$6$}}{\tiny{$7$}}{\tiny{$8$}}{\tiny{$9$}}{\tiny{$10$}}} \\
                \bitboxes{2}{{}{$A$}{$B$}{$C$}{$D$}{$E$}{}{$F$}{}{}{$G$}}
            \end{bytefield}
            \caption{物理结构}
            \label{subfig:demo_physical_structure_binary_tree_sequence_storage}
        \end{subfigure}
        \caption{二叉树的顺序存储示例}
        \label{fig:demo_sequence_storage_binary_tree}
    \end{figure}
\end{frame}

\begin{frame}
    \frametitle{\insertsubsectionhead}
    \begin{figure}
        \centering
        \begin{subfigure}[b]{0.35\textwidth}
            \centering
            \begin{tikzpicture}[ >=stealth, thick, black!50, %
                    list item/.style={draw=gray, circle, thick}, %
                    level distance=4ex, %
                    level 1/.style={sibling distance=16ex}, %
                    level 2/.style={sibling distance=8ex}, %
                    level 3/.style={sibling distance=4ex}, %
                ]
                \node [terminal] {$A$}
                    child [missing]
                    child {node [terminal] {$B$}
                        child [missing]
                        child {node [terminal] {$C$}
                            child [missing]
                            child {node [terminal] {$D$}}}
                    };
            \end{tikzpicture}
            \caption{逻辑结构}
            \label{subfig:demo_logistic_structure_slanted_tree_sequence_storage}
        \end{subfigure}
        ~
        \begin{subfigure}[b]{0.55\textwidth}
            \centering
            \begin{bytefield}{2}
                % \foreach \i in {0,...,10} {\bitbox[]{2}{\tiny{\i}}} \\
                \bitboxes[]{2}{{\tiny{$0$}}{\tiny{$1$}}{\tiny{$2$}}{\tiny{$3$}}{\tiny{$4$}}{\tiny{$5$}}{\tiny{$6$}}{\tiny{$7$}}{\tiny{$8$}}{\tiny{$9$}}{\tiny{$10$}}{\tiny{$11$}}{\tiny{$12$}}{\tiny{$13$}}{\tiny{$14$}}{\tiny{$15$}}} \\
                \bitboxes{2}{{}{$A$}{}{$B$}{}{}{}{$C$}{}{}{}{}{}{}{}{$D$}}
            \end{bytefield}
            \caption{物理结构}
            \label{subfig:demo_physical_structure_slanted_tree_sequence_storage}
        \end{subfigure}
        \caption{右斜树的顺序存储示例}
        \label{fig:demo_sequence_storage_slanted_tree}
    \end{figure}
\end{frame}

\begin{frame}
    \frametitle{\insertsubsectionhead}
    \begin{exampleblock}{顺序存储的特点}
        \begin{itemize}
            \item 可利用性质戊快速访问各结点：$O(1)$
            \item 增删结点可能需要大幅调整存储
            \item 在存储含有稀疏结点的二叉树时需耗费大量存储空间
            \item 仅适合存储含有稠密结点的完全二叉树
        \end{itemize}
    \end{exampleblock}
\end{frame}

\begin{fragile}
    \frametitle{\insertsubsectionhead}
    \bicolumns[0.5]{
        \begin{block}{链式存储}
            \begin{itemize}
                \item 采用二/三叉链表表示结点
                \item 结点除存信息包括
                    \begin{itemize}
                        \item 数据信息
                        \item 左、右子结点地址
                        \item 可选的父结点地址
                    \end{itemize}
                \item 为后续处理方便，设置虚拟首结点
            \end{itemize}
        \end{block}
        \begin{figure}
            \centering
            \begin{bytefield}[boxformatting=\baselinealign]{2}
                \bitbox{3}{\mintinline[fontsize=\smaller]{c}{data}}
                \bitbox{3}{\mintinline[fontsize=\smaller]{c}{left}}
                \bitbox{4}{\mintinline[fontsize=\smaller]{c}{right}}
                \bitbox{4}{\mintinline[fontsize=\smaller]{console}{parent}}
            \end{bytefield}
            \caption{二叉树结点的链式存储结构}
            \label{fig:linked_node_binary_tree}
        \end{figure}
    }{
        \begin{minted}{c}
            typedef struct BinaryTreeNode {
                DataType data; // 数据信息
                struct BinaryTreeNode *left; // 左子结点地址
                struct BinaryTreeNode *right; // 右子结点地址
                struct BinaryTreeNode *parent; // 可选的父结点地址
            } BinaryTreeNode;
        \end{minted}
        \begin{minted}{c}
            typedef struct {
                BinaryTreeNode *head; // 虚拟首结点
            } BinaryTree;
        \end{minted}
    }
\end{fragile}

\begin{fragile}
    \frametitle{\insertsubsectionhead}
    \begin{figure}
        \centering
        \begin{subfigure}[b]{0.3\textwidth}
            \centering
            \begin{tikzpicture}[ >=stealth, thick, black!50, %
                    list item/.style={draw=gray, circle, thick}, %
                    level distance=4ex, %
                    level/.style={sibling distance=12ex/#1}
                ]
                \node [terminal] {$A$}
                    child {node [terminal] {$B$}
                        child {node [terminal] {$D$}}
                        child {node [terminal] {$E$}
                            child {node [terminal] {$G$}}
                            child [missing]
                        }
                    }
                    child {node [terminal] {$C$}
                        child [missing]
                        child {node [terminal] {$F$}}
                    };
            \end{tikzpicture}
            \caption{逻辑结构}
            \label{subfig:demo_logistic_structure_binary_tree_linked_storage}
        \end{subfigure}
        ~
        \begin{subfigure}[b]{0.6\textwidth}
            \centering
            \begin{bytefield}{2}
                \bitbox[]{2}{\tikzmark{tchu}} \\
                \bitbox{2}{首}\bitboxes{1}{{\tikzmark{scha}}{\tikzmark{schx}}{\tikzmark{schy}}}
                \bitbox[]{9}{}\bitbox[]{1}{\tikzmark{tcal}}
                \bitbox{2}{$A$}\bitboxes{1}{{\tikzmark{scab}}{\tikzmark{scac}}{\tikzmark{scah}}} \\
                \bitbox[]{5}{}
                \bitbox[]{2}{\tikzmark{tcbu}}\bitbox[]{8}{}
                \bitbox[]{2}{\tikzmark{tcad}}\bitbox[]{8}{}
                \bitbox[]{2}{\tikzmark{tccu}} \\
                \bitbox[]{5}{}
                \bitbox{2}{$B$}\bitboxes{1}{{\tikzmark{scbd}}{\tikzmark{scbe}}{\tikzmark{scba}}}
                \bitbox[]{15}{}
                \bitbox{2}{$C$}\bitboxes{1}{{\tikzmark{sccx}}{\tikzmark{sccf}}{\tikzmark{scca}}} \\
                \bitbox[]{2}{\tikzmark{tcdu}}\bitbox[]{3}{}\bitbox[]{2}{\tikzmark{tcbd}}\bitbox[]{3}{}
                \bitbox[]{2}{\tikzmark{tceu}}\bitbox[]{13}{}\bitbox[]{2}{\tikzmark{tccd}}\bitbox[]{3}{}
                \bitbox[]{2}{\tikzmark{tcfu}} \\
                \bitbox{2}{$D$}\bitboxes{1}{{\tikzmark{scdx}}{\tikzmark{scdy}}{\tikzmark{scdb}}}
                \bitbox[]{5}{}
                \bitbox{2}{$E$}\bitboxes{1}{{\tikzmark{sceg}}{\tikzmark{scex}}{\tikzmark{sceb}}}
                \bitbox[]{15}{}
                \bitbox{2}{$F$}\bitboxes{1}{{\tikzmark{scfx}}{\tikzmark{scfy}}{\tikzmark{scfc}}} \\
                \bitbox[]{8}{}
                \bitbox[]{2}{\tikzmark{tcgu}}\bitbox[]{2}{\tikzmark{tced}} \\
                \bitbox[]{8}{}
                \bitbox{2}{$G$}\bitboxes{1}{{\tikzmark{scgx}}{\tikzmark{scgy}}{\tikzmark{scge}}}
            \end{bytefield}
            \begin{tikzpicture}[ >=stealth, thick, black!50, overlay, remember picture, %
                    list item/.style={draw=gray, rounded corners, thick} %
                ]
                \draw[->] ($(pic cs:scha)+(0,-0.03)$) to [out=90,in=180] ($(pic cs:tcal)+(0.1,-0.03)$);
                \fill ($(pic cs:scha)+(0,-0.03)$) circle [radius=0.05];
                \draw[->] ($(pic cs:scab)+(0,-0.03)$) to [out=-90,in=90] ($(pic cs:tcbu)+(0,-0.3)$);
                \fill ($(pic cs:scab)+(0,-0.03)$) circle [radius=0.05];
                \draw[->] ($(pic cs:scac)+(0,-0.03)$) to [out=-90,in=90] ($(pic cs:tccu)+(0,-0.3)$);
                \fill ($(pic cs:scac)+(0,-0.03)$) circle [radius=0.05];
                \draw[->] ($(pic cs:scbd)+(0,-0.03)$) to [out=-90,in=90] ($(pic cs:tcdu)+(0,-0.3)$);
                \fill ($(pic cs:scbd)+(0,-0.03)$) circle [radius=0.05];
                \draw[->] ($(pic cs:scbe)+(0,-0.03)$) to [out=-90,in=90] ($(pic cs:tceu)+(0,-0.3)$);
                \fill ($(pic cs:scbe)+(0,-0.03)$) circle [radius=0.05];
                \draw[->] ($(pic cs:sccf)+(0,-0.03)$) to [out=-90,in=90] ($(pic cs:tcfu)+(0,-0.3)$);
                \fill ($(pic cs:sccf)+(0,-0.03)$) circle [radius=0.05];
                \draw[->] ($(pic cs:sceg)+(0,-0.03)$) to [out=-90,in=90] ($(pic cs:tcgu)+(0,-0.3)$);
                \fill ($(pic cs:sceg)+(0,-0.03)$) circle [radius=0.05];
                \draw[dashed,->] ($(pic cs:scah)+(0,-0.03)$) to [out=90,in=90] ($(pic cs:tchu)+(0,-0.3)$);
                \fill ($(pic cs:scah)+(0,-0.03)$) circle [radius=0.05];
                \draw[dashed,->] ($(pic cs:scba)+(0,-0.03)$) to [out=90,in=-90] ($(pic cs:tcad)+(0,0.23)$);
                \fill ($(pic cs:scba)+(0,-0.03)$) circle [radius=0.05];
                \draw[dashed,->] ($(pic cs:scca)+(0,-0.03)$) to [out=90,in=-90] ($(pic cs:tcad)+(0,0.23)$);
                \fill ($(pic cs:scca)+(0,-0.03)$) circle [radius=0.05];
                \draw[dashed,->] ($(pic cs:scdb)+(0,-0.03)$) to [out=90,in=-90] ($(pic cs:tcbd)+(0,0.23)$);
                \fill ($(pic cs:scdb)+(0,-0.03)$) circle [radius=0.05];
                \draw[dashed,->] ($(pic cs:sceb)+(0,-0.03)$) to [out=90,in=-90] ($(pic cs:tcbd)+(0,0.23)$);
                \fill ($(pic cs:sceb)+(0,-0.03)$) circle [radius=0.05];
                \draw[dashed,->] ($(pic cs:scfc)+(0,-0.03)$) to [out=90,in=-90] ($(pic cs:tccd)+(0,0.23)$);
                \fill ($(pic cs:scfc)+(0,-0.03)$) circle [radius=0.05];
                \draw[dashed,->] ($(pic cs:scge)+(0,-0.03)$) to [out=90,in=-90] ($(pic cs:tced)+(0,0.23)$);
                \fill ($(pic cs:scge)+(0,-0.03)$) circle [radius=0.05];
                \draw ($(pic cs:schx)+(-0.1,-0.3)$) to ($(pic cs:schx)+(0.1,0.23)$);
                \draw ($(pic cs:schy)+(-0.1,-0.3)$) to ($(pic cs:schy)+(0.1,0.23)$);
                \draw ($(pic cs:sccx)+(-0.1,-0.3)$) to ($(pic cs:sccx)+(0.1,0.23)$);
                \draw ($(pic cs:scdx)+(-0.1,-0.3)$) to ($(pic cs:scdx)+(0.1,0.23)$);
                \draw ($(pic cs:scdy)+(-0.1,-0.3)$) to ($(pic cs:scdy)+(0.1,0.23)$);
                \draw ($(pic cs:scex)+(-0.1,-0.3)$) to ($(pic cs:scex)+(0.1,0.23)$);
                \draw ($(pic cs:scfx)+(-0.1,-0.3)$) to ($(pic cs:scfx)+(0.1,0.23)$);
                \draw ($(pic cs:scfy)+(-0.1,-0.3)$) to ($(pic cs:scfy)+(0.1,0.23)$);
                \draw ($(pic cs:scgx)+(-0.1,-0.3)$) to ($(pic cs:scgx)+(0.1,0.23)$);
                \draw ($(pic cs:scgy)+(-0.1,-0.3)$) to ($(pic cs:scgy)+(0.1,0.23)$);
            \end{tikzpicture}
            \vspace{-2ex}
            \caption{物理结构}
            \label{subfig:demo_physical_structure_binary_tree_linked_storage}
        \end{subfigure}
        \caption{二叉树的链式存储示例}
        \label{fig:demo_linked_storage_binary_tree}
    \end{figure}
\end{fragile}
\section{二叉搜索树}

\begin{fragile}
    \frametitle{\insertsectionhead}
    \bicolumns[0.55]{
        \begin{block}{二叉搜索树(Binary Search Tree, BST)}
            \begin{itemize}
                \item 对任一结点，左子树的键均小于该结点的键，右子树的键均大于该结点的键
                \item 左右子树也分别是二叉搜索树；这里假设键不重复
                \item 中序遍历得到键的升序序列
            \end{itemize}
        \end{block}
    }{
        \begin{figure}
            \centering
            \begin{tikzpicture}[level distance=6ex, level/.style={sibling distance=14ex/#1}]
                \node[terminal] {$21$}
                    child {node[terminal] {$10$}
                        child {node[terminal] {$5$}}
                        child {node[terminal] {$18$}}}
                    child {node[terminal] {$25$}
                        child {node[terminal] {$22$}}
                        child {node[terminal] {$30$}}};
            \end{tikzpicture}
            \caption{二叉搜索树示例}
        \end{figure}
    }
\end{fragile}

\begin{fragile}
    \frametitle{\insertsectionhead}
    \begin{block}{查找与插入}
        \begin{itemize}
            \item 从根出发比较目标键：较小走左子树，较大走右子树，相等则找到
            \item 到达空指针仍未找到时查找失败；插入可在此位置创建新结点
            \item 本例不插入重复键；也可另行约定计数或重复键处理方式
        \end{itemize}
    \end{block}
    \begin{algorithmx}[alg:bst_search]{二叉搜索树查找}
        \Procedure{Search}{$p,key$}
        \While{$p\neq\mathrm{NULL}$ \textbf{and} $p.key\neq key$}
        \If{$key<p.key$}
        \State $p\gets p.left$
        \Else
        \State $p\gets p.right$
        \EndIf
        \EndWhile
        \State \textbf{return} $p$
        \EndProcedure
    \end{algorithmx}
\end{fragile}

\begin{frame}
    \frametitle{\insertsectionhead}
    \begin{block}{删除结点}
        \begin{enumerate}
            \item 无孩子：直接删除该叶结点
            \item 只有一个孩子：用孩子替代该结点
            \item 有两个孩子：用右子树中最小结点（中序后继）的键替换，
                  再删除该后继结点；后继至多有一个孩子
        \end{enumerate}
    \end{block}
    \begin{exampleblock}{例：删除上一图中的$21$}
        中序后继为$22$；将根键替换为$22$，再删除原位置的叶结点$22$。
        结果仍满足左小右大的性质。
    \end{exampleblock}
\end{frame}

\begin{frame}
    \frametitle{\insertsectionhead}
    \begin{block}{性能与取舍}
        \begin{itemize}
            \item 查找、插入和删除沿树高$h$进行，时间均为$O(h)$
            \item 高度平衡时$h=O(\log n)$；退化成链时$h=O(n)$
            \item 顺序插入$1,2,\ldots,n$可能产生退化树；平衡搜索树可控制高度
            \item BST支持按序遍历和范围查询；普通散列表不保持键的顺序
        \end{itemize}
    \end{block}
    \begin{alertblock}{思考}
        依次插入$21,10,25,18$后，中序遍历结果是什么？
    \end{alertblock}
    \begin{block}<2-|handout:0>{参考答案}
        $10,18,21,25$。
    \end{block}
\end{frame}


\section{小结}

\begin{frame}
    \frametitle{\insertsectionhead}
    \begin{itemize}
        \item 半线性结构以树与森林为代表，前驱为线性、后继为非线性
        \item 树与森林的问题均可转化为二叉树，使研究更加方便、统一
        \item 二叉树可采用顺序与链式两种方式存储
        \item 二叉树的遍历可将半线性结构转化为线性结构
        \item 二叉搜索树利用键的顺序进行查找；操作代价取决于树高
        \item Huffman编码树可用于数据压缩
    \end{itemize}
\end{frame}

\begin{frame}{思考与自测}
    \begin{itemize}
        \item 根结点为$A$，左孩子为$B$，右孩子为$C$：写出三种深度优先遍历。
        \item 若树有$n$个结点，遍历每个结点一次的时间复杂度是多少？
    \end{itemize}
    \pause
    \begin{block}<2-|handout:0>{参考答案}
        前序$A,B,C$；中序$B,A,C$；后序$B,C,A$。遍历时间为$O(n)$。
    \end{block}
\end{frame}

\begin{standout}[]
    问与答
\end{standout}